% Part: proof-theory
% Chapter: proof-search
% Section: completeness

\documentclass[../../../include/open-logic-section]{subfiles}

\begin{document}

\olfileid{pt}{ps}{cpl}

\olsection{પૂર્ણતા}

હવે સાબિતી-શોધ
અલ્ગોરિધમ અચૂક
છે તે બતાવીએ:
જો તે અધવચ્ચે
અટકી જાય અથવા
કદી પૂર્ણ ન
થાય, તો શરૂઆતના
સિક્વન્ટ
$\Gamma \Sequent \Delta$ની
કોઈ !!{proof}
નથી. તે માટે
!!a{structure}~$\Struct M$
એવી બનાવીએ કે
દરેક~$!A \in \Gamma$
માટે~$\Sat{M}{!A}$
અને દરેક
$!A \in \Delta$
માટે~$\Sat/{M}{!A}$.
બીજા શબ્દોમાં,
$\Gamma$નાં બધાં
!!{formula}s
અને~$\Delta$નાં
બધાં !!{formula}s
અનુક્રમે
$\Struct M$માં
સત્ય અને અસત્ય
છે. એટલે
$\Gamma \Sequent \Delta$
પ્રામાણ્ય નથી.
\Log{G3c} યથાર્થ
હોવાથી
$\Gamma \Sequent \Delta$ની
!!{proof} નથી.
\sourcecorrection{OLMD-269}{મૂળમાં $\Delta$ના દરેક સભ્ય માટેની શરતમાં $! \in \Delta$ લખાયું છે; અહીં સૂત્રચલ $!A \in \Delta$ હોવો જોઈએ.}

$\Struct{M}$ને
!!{formula}sની
જોડી~$\Theta, \Xi$
અને
!!{constant}sના
ગણ~$C$માંથી
વ્યાખ્યાયિત કરીએ.
આ~$\Theta, \Xi$
અને~$C$
સાબિતી-શોધનો
અસફળ (અટકેલી
અથવા અપૂર્ણ
રહેલી) દોડની
\emph{નિષ્ફળતા-શાખા}
માંથી મળે છે.
નિષ્ફળતા-શાખા
સિક્વન્ટોનો
ક્રમ~$\Pi_n \Sequent
\Lambda_n$
અને અનુરૂપ
!!{constant}sના
ગણ~$C_n$નો
બને છે, જ્યાં
$\Pi_0 \Sequent \Lambda_0$
એ અંતિમ
સિક્વન્ટ
$\Gamma \Sequent
\Delta$ છે,
અને
$\Pi_{n+1} \Sequent \Lambda_{n+1}$
એ~$\Pi_n
\Sequent \Lambda_n$નું
તબક્કા~$n+1$એ
બનેલું એવું
આધારવિધાન છે
જેની !!a{proof}
અલ્ગોરિધમ શોધતું
નથી. શાખા
ચાલુ હોય તેવા
દરેક~$n$
માટે આવું
આધારવિધાન હોય
જ છે, નહીં તો
અલ્ગોરિધમને
!!a{proof}
મળી ગઈ હોત.
$C_n$ એ
$\Pi_n
\Sequent \Lambda_n$
સાથે સંકળાયેલો
!!{constant}sનો
ગણ છે.
\sourcecorrection{OLMD-270}{મૂળનું છેલ્લું સિક્વન્ટ $\Pi_n \Sequent \Delta_n$ લખાયું છે; નિષ્ફળતા-શાખાનો જમણો ભાગ સમગ્ર વ્યાખ્યામાં $\Lambda_n$ છે.}
\sourcecorrection{OLMD-271}{મૂળ કહે છે કે દરેક $n$ માટે આગળનું આધારવિધાન હોવું જ જોઈએ; સીમિત નિષ્ફળતા-શાખા આણ્વિક સિક્વન્ટે પૂરી થાય છે. આગળનું આધારવિધાન ફક્ત શાખા ચાલુ હોય એવા તબક્કે જરૂરી છે.}

અલ્ગોરિધમ માત્ર
આણ્વિક
!!{formula}sવાળા
સૌથી ઉપરના
સિક્વન્ટ પર
અટકી જાય,
તો ત્યાં પૂરી
થતી શાખા
નિષ્ફળતા-શાખા
છે. જો
અલ્ગોરિધમ
કદી પૂર્ણ ન
થાય, તો તે
વધુ ને વધુ
મોટાં વૃક્ષો
બનાવતું રહે.
આને અનંત
વૃક્ષ વધતું
હોય તેમ
વિચારી શકાય
(અલ્ગોરિધમનાં
અનંત તબક્કા
પછીનું વૃક્ષ).
આ અનંત વૃક્ષ
સીમિત શાખાવાળું
છે: દરેક ગાંઠને
માત્ર એક કે
બે સંતાન
સિક્વન્ટ છે,
જે તેની તરત
ઉપર છે.
K\H{o}nigના
ઉપપાદ્ય પ્રમાણે
દરેક અનંત,
સીમિત શાખાવાળા
વૃક્ષમાં અનંત
શાખા હોય છે.
શોધ અલ્ગોરિધમ
કાયમ ચાલે
ત્યારે આવી
અનંત શાખા
નિષ્ફળતા-શાખા
છે.

$\Pi_n \Sequent \Lambda_n$
અને~$C_n$નો
ક્રમ નિષ્ફળતા-શાખા
હોય, તો
$\Theta = \bigcup_n \Pi_n$
અને
$\Xi = \bigcup_n \Lambda_n$
લઈએ
(!!{formula}sના
અનુક્રમણિકા અને
પુનરાવર્તિત
આવૃત્તિઓ દૂર
કરીને), અને
$C =
\bigcup_n C_n$
લઈએ.

હવે~$\Struct{M}$
વ્યાખ્યાયિત કરીએ.
\begin{enumerate}
\item પ્રદિશ
$\Domain{M}$ એ
$C$ અને
$\Gamma \Sequent \Delta$માં
આવતાં
!!{function}s
હોય તો તેમની
મદદથી, પણ
!!{variable}s વિના,
બનતા પદોનો
ગણ છે.
\item $c \in \Domain{M}$
હોય તો
$\Assign{c}{M} = c$.
\item $f$ એ
$\Gamma \Sequent \Delta$માંનું
$m$-સ્થાની
!!{function}
હોય, અને
$t_1$,
\dots,
$t_m \in \Domain{M}$
હોય, તો
$\Assign{f}{M}(t_1, \dots, t_m) =
f(t_1, \dots, t_m)$.
\sourcecorrection{OLMD-272}{મૂળમાં $m$-સ્થાની વિધેયના મૂલ્ય માટે છેલ્લે $f(t_1,\dots,t_n)$ લખાયું છે; ઇનપુટ અને વિધેયની સ્થાનસંખ્યા મુજબ અંતિમ પદ $t_m$ છે.}
\item $R$ એ
$\Gamma \Sequent
\Delta$માંનું
$m$-સ્થાની
!!{predicate}
હોય, તો
$\Assign{R}{M} = \Setabs{\tuple{t_1, \dots, t_m} \in
\Domain{M}^m}{R(t_1, \dots, t_m) \in \Theta}$.
\sourcecorrection{OLMD-273}{મૂળમાં $m$-સ્થાની સંબંધ માટે પ્રદિશની ઘાત અને $R$ના અંતિમ પદમાં $n$ છે; બંને જગ્યાએ સ્થાનસંખ્યા $m$ જ હોવી જોઈએ.}
\end{enumerate}
$\Struct{M}$ એક
\emph{પદ-નિદર્શ}
છે: તેનો પ્રદિશ
પદોનો ગણ છે,
અને
!!{constant}s તથા
!!{function}sનું
અર્થઘટન એવી
રીતે થાય છે
કે
$\Value{t}{M} = t$.
$\Theta$નાં
આણ્વિક
!!{formula}s
વડે
$\Assign{R}{M}$
એવી રીતે
વ્યાખ્યાયિત
થાય છે કે
$\Sat{M}{R(t_1, \dots, t_n)}$
ત્યારે અને
ત્યારે જ,
જ્યારે
$R(t_1, \dots, t_n)
\in \Theta$.

\begin{lem}\ollabel{lem:truth}
દરેક
$!A \in \Theta \cup \Xi$
માટે:
\begin{enumerate}
  \item $!A \in \Theta$
  હોય તો
  $\Sat{M}{!A}$.
  \item $!A \in \Xi$
  હોય તો
  $\Sat/{M}{!A}$.
\end{enumerate}
\end{lem}

\begin{proof}
  $\depth{!A}$
  પર આગમન
  કરીએ.

પહેલાં માનો કે
$!A$ આણ્વિક છે,
એટલે કાં તો
$!A \ident \lfalse$
અથવા
$!A
\ident R(t_1, \dots, t_m)$.

$\Pi_n \Sequent \Lambda_n$
નિષ્ફળતા-શાખા
હોવાથી કોઈ
$\Pi_n$માં
$\lfalse$ નથી
(નહિતર
$\Pi_n \Sequent \Lambda_n$
સ્વયંસિદ્ધ
થાય).
$\Struct{M}$ની
વ્યાખ્યા પ્રમાણે
$\Sat{M}{R(t_1, \dots, t_m)}$
ત્યારે અને
ત્યારે જ,
જ્યારે
$R(t_1, \dots, t_m) \in
\Theta$.
બંને વાતો
મળીને આપે છે:
$!A \in
\Theta$ હોય,
તો
$\Sat{M}{!A}$.

$!A$ આણ્વિક
હોય અને
$!A \in \Pi_n$
હોય, તો
$!A \in \Pi_{n+1}$;
તેવી જ રીતે
$!A \in \Lambda_n$
હોય તો
$!A \in \Lambda_{n+1}$.
(સાબિતી-શોધ
અલ્ગોરિધમ
ઉત્તરાધિકારી
સિક્વન્ટો જે
રીતે બનાવે
છે તેમાંથી
આ ફલિત થાય
છે.) તેથી
$!A \in \Theta \cap \Xi$
હોય, તો કોઈ
$n$ માટે
$!A \in \Pi_n \cap
\Lambda_n$.
એટલે
$\Pi_n \Sequent \Lambda_n$
સ્વયંસિદ્ધ
થાય, પરંતુ
નિષ્ફળતા-શાખામાં
સ્વયંસિદ્ધ ન
હોઈ શકે.
આથી આણ્વિક
$!A$ માટે
$!A \notin \Theta
\cap \Xi$,
અને પરિણામે
$!A \in \Xi$
હોય તો
$!A \notin \Theta$.
$\Struct{M}$ની
વ્યાખ્યાથી
આનો અર્થ
એ કે
$!A \in \Xi$
હોય, તો
$\Sat/{M}{!A}$.

આગમન-પગલામાં
માનો કે~$!A$
કરતાં ઓછી
!!{depth} ધરાવતા
બધા !!{formula}s
માટે (1)
અને (2)
સાચાં છે.
$!A$ માટે
(1) અને
(2) કિસ્સા
જુદા પાડીને
સાબિત કરીએ.

પહેલાં નોંધો કે
$!A^i$
સિક્વન્ટ
$\Pi_n \Sequent \Lambda_n$માં
હોય, તો~$i$
તે જ
સિક્વન્ટ~$\Pi_n \Sequent \Lambda_n$નો
સૌથી
નાનો અનુક્રમણિકાંક
ન હોય ત્યાં
સુધી~$!A^i$
પછીના
$\Pi_{n+1} \Sequent \Lambda_{n+1}$
માં પણ રહે
છે. દરેક તબક્કે
ઉમેરાતા સૌથી
ઉપરના સિક્વન્ટોમાં
સૌથી નાનો
અનુક્રમણિકાંક
વધે છે.
તેથી આખરે
એવું
$\Pi_n \Sequent \Lambda_n$
આવે છે જેમાં
$!A^i$ હોય
અને~$i$
સૌથી નાનો
હોય. તબક્કા
$n+1$એ
અલ્ગોરિધમ~$!A^i$ને
ઘટાડે છે.
એટલે
$!A \in \Theta$
હોય, તો કોઈ
$n$ માટે
$\Pi_n = \Pi_n',
!A^i$
અને~$i$
સૌથી નાનો
અનુક્રમણિકાંક
હોય. તેવી જ
રીતે
$!A \in \Xi$
હોય, તો કોઈ
$n$ માટે
$\Lambda_n = \Lambda_n', !A^i$
અને~$i$
સૌથી નાનો
અનુક્રમણિકાંક
હોય.

\begin{enumerate}
  \item $!A \in \Theta$
  અને
  $!A \ident !B \land !C$
  હોય. ત્યારે
  કોઈ~$n$ માટે
  $\Pi_n = \Pi_n', (!B \land !C)^i$.
  $\Pi_{n+1} \Sequent
  \Lambda_{n+1}$
  એ~\LeftR{\land}નું
  અનુરૂપ
  આધારવિધાન
  છે, એટલે
  $\Pi_{n+1} = \Pi_n', !B^k, !C^{k+1}$.
  તેથી
  $!B, !C \in
  \Theta$.
  આગમન-પરિકલ્પનાથી
  $\Sat{M}{!B}$
  અને
  $\Sat{M}{!C}$.
  $\Sat{M}{}$ની
  વ્યાખ્યા પ્રમાણે
  $\Sat{M}{!B \land !C}$.

  \item $!A \in \Xi$
  અને
  $!A \ident !B \land !C$
  હોય. ત્યારે
  કોઈ~$n$ માટે
  $\Lambda_n = \Lambda_n', (!B \land !C)^i$.
  $\Pi_{n+1} \Sequent
  \Lambda_{n+1}$
  એ નિષ્કર્ષ
  $\Pi_n \Sequent \Lambda'_n, (!B \land !C)^i$
  ધરાવતા
  \RightR{\land}ના
  બે આધારવિધાનોમાંથી
  એક છે, એટલે
  $\Lambda_{n+1} = \Lambda_n', !B^k$
  અથવા
  $\Lambda_{n+1} = \Lambda_n',
  !C^k$.
  તેથી
  $!B \in \Xi$
  અથવા
  $!C \in \Xi$.
  આગમન-પરિકલ્પનાથી
  $\Sat/{M}{!B}$
  અથવા
  $\Sat/{M}{!C}$.
  $\Sat{M}{}$ની
  વ્યાખ્યાથી
  $\Sat/{M}{!B \land !C}$.
  \sourcecorrection{OLMD-274}{મૂળના જમણા સંયોજન-કિસ્સામાં ઘટાડાતા મુખ્ય સૂત્રના બે ઉલ્લેખમાં અનુક્રમણિકાંક $i$ છૂટ્યો છે; અલ્ગોરિધમ અનુક્રમણિકાંકવાળા સિક્વન્ટો પર ચાલે છે.}
  
  \item
  $!A \ident \lnot !B$,
  $!A \ident !B \lor !C$
  અને
  $!A \ident !B \lif !C$
  કિસ્સા સમાન
  છે અને
  કસરત તરીકે
  મૂક્યાં છે.

  \item $!A \in \Theta$
  અને
  $!A \ident \lexists[x][!B(x)]$
  હોય. ત્યારે
  કોઈ~$n$ માટે
  $\Pi_n = \Pi_n', \lexists[x][!B(x)]^i$.
  $\Pi_{n+1}
  \Sequent \Lambda_{n+1}$
  એ~\LeftR{\lexists}નું
  અનુરૂપ
  આધારવિધાન
  છે, એટલે
  કોઈ~$c$ માટે
  $\Pi_{n+1} = \Pi_n', !B(c)^k$.
  તેથી
  $!B(c) \in \Theta$
  અને
  $c \in C$.
  આગમન-પરિકલ્પનાથી
  $\Sat{M}{!B(c)}$.
  $\Sat{M}{}$ની
  વ્યાખ્યાથી
  $\Sat{M}{\lexists[x][!B(x)]}$.

  \item $!A \in \Xi$
  અને
  $!A \ident \lexists[x][!B(x)]$
  હોય. ત્યારે
  કોઈ~$n$ માટે
  $\Lambda_n = \Lambda_n', \lexists[x][!B(x)]^i$.
  જ્યારે જ્યારે
  $\lexists[x][!B(x)]$
  ઘટાડાય છે,
  ત્યારે
  $\lexists[x][!B(x)]$
  નવા અનુક્રમણિકાંક
  સાથે આધાર
  સિક્વન્ટની
  જમણી બાજુએ
  રહે છે.
  તેથી નિષ્ફળતા-શાખામાં
  તે જમણી બાજુ
  આવે તો
  $\lexists[x][!B(x)]$
  અનંત વાર
  ઘટાડાય છે.
  દરેક
  $t \in
  \Domain{M}$
  ક્યારેક
  નિષ્ફળતા-શાખામાં
  ``ફક્ત~$C_n$નાં
  અચળો ધરાવતું,
  જમણી બાજુના
  $\lexists[x][!B(x)]$ના
  ઘટાડામાં હજી
  ન વપરાયેલું
  પ્રથમ પદ''
  બને છે.
  તેથી દરેક
  $t \in \Domain{M}$
  માટે કોઈ~$n$
  પર
  $!B(t) \in \Lambda_n$,
  અને પરિણામે
  $!B(t) \in \Xi$.
  આગમન-પરિકલ્પનાથી
  દરેક
  $t \in
  \Domain{M}$
  માટે
  $\Sat/{M}{!B(t)}$.
  $\Sat{M}{}$ની
  વ્યાખ્યાથી
  $\Sat/{M}{\lexists[x][!B(x)]}$.

  \item
  $!A \ident \lforall[x][!B(x)]$
  કિસ્સા સમાન
  છે અને
  કસરત તરીકે
  મૂક્યાં છે.
\end{enumerate}
\end{proof}

\begin{cor}\ollabel{cor:complete}
\Log{G3c}નો
સાબિતી-શોધ
અલ્ગોરિધમ પૂર્ણ
છે: તે અધવચ્ચે
અટકે અથવા
કદી પૂર્ણ ન
થાય, તો શરૂઆતનો
સિક્વન્ટ
$\Gamma \Sequent \Delta$
પ્રામાણ્ય નથી
અને તેથી
તેની !!{proof}
નથી.
\end{cor}

\begin{proof}
  અલ્ગોરિધમ
  અટકે અથવા
  કદી પૂર્ણ ન
  થાય, તો
  નિષ્ફળતા-શાખા
  મળે છે;
  તેમાંથી
  $\Struct{M}$
  વ્યાખ્યાયિત
  કરી શકાય.
  \olref{lem:truth}
  મુજબ દરેક
  $!A \in \Gamma$
  માટે
  $\Sat{M}{!A}$
  અને દરેક
  $!A \in \Delta$
  માટે
  $\Sat/{M}{!A}$.
  (યાદ રાખો:
  $\Pi_0 = \Gamma$
  અને
  $\Lambda_0 =
  \Delta$,
  તેથી
  $\Gamma \subseteq \Theta$
  અને
  $\Delta \subseteq \Xi$.)
  એટલે
  $\Gamma \Sequent \Delta$
  પ્રામાણ્ય નથી.
  \Log{G3c}ની
  યથાર્થતા મુજબ
  $\Gamma \Sequent \Delta$ની
  !!{proof}
  નથી.
\end{proof}

\begin{cor}
\Log{G3c}
પૂર્ણ છે: જો
$\Gamma \Sequent \Delta$
પ્રામાણ્ય હોય, તો
$\Log{G3c} \Proves \Gamma \Sequent \Delta$.
\end{cor}

\begin{proof}
  પ્રતિવિધાન
  સાબિત કરીએ.
  માનો કે
  $\Log{G3c} \Proves/ \Gamma
  \Sequent \Delta$.
  શોધવાની
  !!{proof}
  નથી, તેથી
  સાબિતી-શોધ
  અલ્ગોરિધમ
  અટકે અથવા
  કદી પૂર્ણ ન
  થાય. પછી
  \olref{cor:complete}
  મુજબ
  $\Gamma \Sequent \Delta$
  પ્રામાણ્ય નથી.
\end{proof}


\end{document}
